\section{Setup}\label{sec:setup}

Consider predicting labels $y \in \sY$ from input features $x \in \sX$.
Given a model family $\Theta$,
loss $\ell: \Theta \times (\sX \times \sY) \to \R_+$,
and training data drawn from some distribution $P$,
the standard goal is to find a model $\theta \in \Theta$ that minimizes the expected loss $\E_P[\ell(\theta; (x, y)]$ under the same distribution $P$.
The standard training procedure for this goal is empirical risk minimization (ERM):
\begin{align}\label{eqn:hterm}
  \hterm \eqdef \argmin_{\theta \in \Theta} \ \E_{(x, y) \sim \hP}[\ell(\theta; (x, y))],
\end{align}
where $\hP$ is the empirical distribution over the training data.

In distributionally robust optimization (DRO) \citep{bental2013robust, duchi2016}, we aim instead to minimize the worst-case expected loss over
an uncertainty set of distributions $\sQ$:
\begin{align}\label{eq:dro}
  \underset{\theta \in \Theta}{\vphantom{\sup}\min} \Bigl\{ \sR(\theta) \eqdef \sup_{Q \in \sQ} \E_{(x, y) \sim Q}[\ell(\theta; (x, y))] \Bigr\}.
\end{align}
The uncertainty set $\sQ$ encodes the possible test distributions that we want our model to perform well on.
Choosing a general family $\sQ$, such as a divergence ball around the training distribution,
confers robustness to a wide set of distributional shifts,
but can also lead to overly pessimistic models which optimize for implausible worst-case distributions
\citep{duchi2019distributionally}.

To construct a realistic set of possible test distributions without being overly conservative,
we leverage prior knowledge of spurious correlations to define groups over the training data and then define the uncertainty set $\sQ$ in terms of these groups.
Concretely, we adopt the \emph{group DRO} setting \citep{hu2018does,oren2019drolm} where the training distribution $P$ is assumed to be a mixture of $m$ groups $P_g$ indexed by $\sG = \{1, 2, \ldots, m\}$.\footnote{
  In our main experiments, $m=4$ or $6$;
  we also use $m=64$ in our supplemental experiments.
}
We define the uncertainty set $\sQ$ as any mixture of these groups,
i.e., $\sQ := \{ \sum_{g=1}^m q_g P_g: q \in \Delta_m \}$,
where $\Delta_m$ is the $(m-1)$-dimensional probability simplex;
this choice of $\sQ$ allows us to learn models that are robust to group shifts.
Because the optimum of a linear program is attained at a vertex,
the worst-case risk~\eqref{eq:dro} is equivalent to a maximum over the expected loss of each group,
\begin{align}
\label{eqn:groupdroq}
  \sR(\theta)
  &= \underset{g \in \sG}{\max} \ \E_{(x, y) \sim P_g}[\ell(\theta; (x, y))].
\end{align}
We assume that we know which group each training point comes from---i.e., the training data comprises $(x, y, g)$ triplets---though we do not assume we observe $g$ at test time, so the model cannot use $g$ directly.
Instead, we learn a \emph{group DRO} model minimizing the empirical worst-group risk $\hat\sR(\theta)$:
\begin{align}\label{eqn:htdro}
  \htdro \eqdef \underset{\theta \in \Theta}{\vphantom{\sup}\argmin} \Bigl\{\hat\sR(\theta) :=  \underset{g \in \sG}{\vphantom{\argmin}\max} \ \E_{(x, y) \sim \hP_g}[\ell(\theta; (x, y))] \Bigr\},
\end{align}
where each group $\hP_g$ is an empirical distribution over all training points $(x, y, g')$ with $g' = g$
(or equivalently, a subset of training examples drawn from $P_g$).
Group DRO learns models with good worst-group \emph{training} loss across groups.
This need not imply good worst-group \emph{test} loss
because of the worst-group \emph{generalization gap}
$\delta \eqdef  \sR(\theta)  -\hat{\sR}(\theta)$.
We will show that for overparameterized neural networks,
$\delta$ is large unless we apply sufficient regularization.

\subsection{Applications}
In the rest of this paper,
we study three applications that share a similar structure (\reffig{dataset}):
each data point $(x, y)$ has
some input attribute $a(x) \in \sA$ that is spuriously correlated with the label $y$,
and we use this prior knowledge to form $m = |\sA| \times |\sY|$ groups,
one for each value of $(a, y)$.
We expect that models that learn the correlation between $a$ and $y$ in the training data
would do poorly on groups for which the correlation does not hold
and hence do worse on the worst-group loss $\sR(\theta)$.

\paragraph{Object recognition with correlated backgrounds (Waterbirds dataset).}
Object recognition models can spuriously rely on the image background instead of learning to recognize the actual object
\citep{ribeiro2016lime}.
We study this by constructing a new dataset, Waterbirds,
which combines bird photographs from the Caltech-UCSD Birds-200-2011 (CUB) dataset \citep{wah2011cub}
with image backgrounds from the Places dataset \citep{zhou2017places}.
We label each bird as one of $\sY = \{\text{waterbird}, \text{landbird}\}$
and place it on one of $\sA = \{\text{water background}, \text{land background}\}$,
with waterbirds (landbirds) more frequently appearing against a water (land) background (\refapp{dataset_details}).
There are $n=4795$ training examples and $56$ in the smallest group
(waterbirds on land).

\paragraph{Object recognition with correlated demographics (CelebA dataset).}
Object recognition models (and other ML models more generally) can also learn spurious associations between the label and
demographic information like gender and ethnicity \citep{buolamwini2018gender}.
We examine this on the CelebA celebrity face dataset \citep{liu2015deep},
using hair color ($\sY = \{\text{blond}, \text{dark}\}$) as the target
and gender ($\sA = \{\text{male}, \text{female}\}$) as the spurious attribute.
There are $n=162770$ training examples in the CelebA dataset, with $1387$ in the smallest group (blond-haired males).

\paragraph{Natural language inference (MultiNLI dataset).}
In natural language inference, the task is to determine if a given hypothesis is entailed by, neutral with, or contradicts a given premise.
Prior work has shown that crowdsourced training datasets for this task
have significant annotation artifacts,
such as the spurious correlation between contradictions and the presence of the negation words \emph{nobody}, \emph{no}, \emph{never}, and \emph{nothing}
\citep{gururangan2018annotation}.
We divide the MultiNLI dataset \citep{williams2018broad}
into $m=6$ groups, one for each pair of
labels $\sY = \{\text{entailed}, \text{neutral}, \text{contradictory}\}$ and
spurious attributes $\sA = \{\text{no negation}, \text{negation}\}$.
There are $n=206175$ examples in our training set,
with $1521$ examples in the smallest group (entailment with negations);
see \refapp{dataset_details} for more details on dataset construction and the training/test split.
